\section{Characters of Finite Fields}\label{sec:chars-finite-fields}
In a finite field $\bbF_q$ there are two inherent group structures--namely, the additive group and the multiplicative group. The nonzero elements of $\bbF_q$ forms a cyclic group of $q-1$ elements under multiplication and the whole $\bbF_q$ forms a group of $q$ elements under addition.

\subsection{Additive Characters of \texorpdfstring{$\bbF_q$}{Fq}}
We first turn to the additive group of $\bbF_q$. Characters of the additive group $\bbF_q$ are called \emph{additive characters}\index{character!additive} of $\bbF_q$. We denote the set of all additive characters of $\bbF_q$ by $\bm{\mcX}_q$. 

Suppose $p$ be the characteristic of $\bbF_q$.  Then the prime field contained in $\bbF_q$ is $\bbF_p$. Let $\tr:\bbF_q\to\bbF_p$ be the absolute \emph{trace function}\index{trace function}\index{trace!of a field extension} from $\bbF_q$ to $\bbF_p$. Then define the function $\chi_1$ by $$\chi_1(\alpha)=e^{2\pi i\frac{\tr(\alpha)}{p}}\quad \forall\ \alpha\in \bbF_q$$Then $\chi_1$ is an additive character of $\bbF_q$ since for all $\alpha_1,\alpha_2\in\bbF_q$, $\tr(\alpha_1+\alpha_2)=\tr(\alpha_1)+\tr(\alpha_2)$ and hence $\chi_1(\alpha_1+\alpha_2)=\chi_1(\alpha_1)\cdot \chi_1(\alpha_2)$. This character $\chi_1$ is called the \emph{canonical additive character}\index{character!canonical additive} of $\bbF_q$.
\begin{Theorem}{Structure of Additive Characters}{thm:add-char-define}
  For $\beta\in\bbF_q$, define $\chi_{\beta}$ with $\chi_{\beta}(\alpha)=\chi_1(\beta\cdot \alpha)$ for all $\alpha\in\bbF_q$. Then $\chi_{\beta}$ is an additive character of $\bbF_q$ and every additive character of $\bbF_q$ is obtained this way i.e. $\forall\ \chi\in\Xq$, there exists $\beta\in\bbF_q$ such that $\forall\ \alpha\in\bbF_q$, $\chi(\alpha)=\chi_{\beta}(\alpha)$. 
\end{Theorem}
\begin{proof}
  Now for any $\alpha_1,\alpha_2\in\bbF_q$, we have $$\chi_{\beta}(\alpha_1+\alpha_2)=\chi_1(\beta(\alpha_1+\alpha_2))=\chi_1(\beta\cdot \alpha_1)\cdot \chi_1(\beta\cdot \alpha_2)=\chi_{\beta}(\alpha_1)\cdot \chi_{\beta}(\alpha_2)$$ Hence $\chi_{\beta}$ is an additive character of $\bbF_q$.

  Now we will show that for any two distinct $\beta_1,\beta_2\in\bbF_q$ with $\beta_1\neq \beta_2$ there exists an $\alpha\in\bbF_q$ such that $\chi_{\beta_1}(\alpha)\neq \chi_{\beta_2}(\alpha)$. Suppose that is not the case. Then for all $\alpha\in\bbF_q$, $$\chi_{\beta_2}(\alpha)=\chi_{\beta_2}(\alpha)\iff \chi_1(\beta_1\cdot \alpha)=\chi_1(\beta_2\cdot \alpha)\iff \chi_1((\beta_1-\beta_2)\alpha)=1.$$Now $\chi_1((\beta_1-\beta_2)\alpha)=1$ for all $\alpha\in \bbF_q$ if and only if $\tr((\beta_1-\beta_2)\alpha)=0$ for all $\alpha\in\bbF_q$. Since  $\tr(X)=X+X^{p}+\cdots+X^{p^{r-1}}$ where $q=p^r$. Since $\tr$ is $\bbF_p-$linear,  $\tr((\beta_1-\beta_2)X)$ has $p^r$ roots in $\bbF_q$ which is more than the degree $p^{r-1}$. Hence $\tr((\beta_1-\beta_2)X)$ is the zero polynomial. Therefore $\beta_1=\beta_2$. So for any two distinct $\beta_1,\beta_2\in\bbF_q$ with $\beta_1\neq \beta_2$ there exists an $\alpha\in\bbF_q$ such that $\chi_{\beta_1}(\alpha)\neq \chi_{\beta_2}(\alpha)$.

  Hence for every $\beta\in\bbF_q$ we have a distinct additive character. Therefore we have $q$ additive characters. Since there are $q$ additive characters of $\bbF_q$ by \thmref{thm:group-char-equal} every additive character of $\bbF_q$ can be obtained the above way.
\end{proof}
 

The orthogonality relations in \obsref{obs:orthogonal-relation} when applied to additive characters and along with \thmref{thm:add-char-define} gives the following result.
\begin{Theorem}{Orthogonality of Additive Characters}{thm:add-char-props}
  \begin{enumerate}[label=(\roman*)]
    \item For any $a\in\bbF_q$, the corresponding additive character is $\chi_a$. Then $$\sum_{\gm\in\bbF_q}\chi_a(\gm)=\begin{cases}
              0 & \text{if $a\neq 0$} \\
              q & \text{if $a=0$}
            \end{cases}$$
    \item For any two additive characters $\chi_{a},\chi_{b}$ where $a,b\in\bbF_q$, $$\sum_{\gm\in\bbF_q}\chi_{a}(\gm)\cdot \ov{\chi_{b}}(\gm)=\begin{cases}
              0 & \text{if $a\neq b$} \\
              q & \text{if $a=b$}
            \end{cases}$$
    \item For any two elements $\alpha,\beta\in\bbF_q$, $$\sum_{c\in\bbF_q}\chi_c(\alpha)\cdot\ov{\chi_c(\beta)}=\begin{cases}
              0 & \text{if $\alpha\neq \beta$} \\
              q & \text{if $\alpha=\beta$}
            \end{cases}$$
  \end{enumerate}
\end{Theorem}

\subsection{Multiplicative Characters of \texorpdfstring{$\bbF_q$}{Fq}}\label{sec:mult-char}
Characters of the multiplicate group $\bbF_q^*$ are called \emph{multiplicative characters}\index{character!multiplicative} of $\bbF_q$. We denote the set of all multiplicative characters of $\bbF_q$ by $\bm{\mcM}_q$. Since $\bbF_q^*$ is a cyclic group of order $q-1$ by \lemref{lem:chars-cyclic-group} we have the following:
\begin{lemma}{Structure of Multiplicative Characters}{lem:mult-char-define}
  The group of multiplicative characters of $\bbF_q$ is a cyclic group of order $q-1$ and are given by $\{\psi_j\colon 0\leq j\leq q-2\}$ where if $g$ is the generator of $\bbF_q^*$ then for all $k\in\{0,1,\dots, q-2\}$ we have $\psi_j(g^k)=e^{2\pi i\frac{jk}{q-1}}$ where $\psi_0$ is the identity element of the group and the trivial character.
\end{lemma}

Every multiplicative character $\psi$ of $\bbF_q$ has the property that $\psi^{q-1}=\psi_0$. Let $\psi$ be the generator of the cyclic group of multiplicative characters. Then $\psi$ is called the \emph{principal character}\index{character!principal}\index{principal character}\index{multiplicative character!exponent $d$}. For any multiplicative character $\psi$ of $\bbF_q$ we say $\psi$ is of \emph{order}\index{character!order} $d$ if $\psi^d=\psi_0$ and $d$ is the smallest positive integer with this property. We use $\ord(\psi)$ to denote the order of $\psi$. Hence $d\mid q-1$. We say $\psi$ is of \emph{exponent}\index{character!exponent} $e$ if $\psi^e=\psi_0$. Clearly $d\mid e$ where $d$ is the order of $\psi$. Let $\Mq^{(e)}$ denote the set of all multiplicative characters of $\Mq$ of exponent $e$. 

\begin{note}
  For odd $q$ a special multiplicative character of interest is the \emph{quadratic character}\index{character!quadratic}\index{quadratic character}\index{quadratic residue}. It is denoted by $\eta$. It only takes the value $\pm1$. So $\ord{\eta}=2$. For all $\alpha\in \bbF_q^*$, $\eta(\alpha)=1$ if $\alpha$ is a \emph{quadratic residue} and $-1$ if $\alpha$ is a non-quadratic residue. Therefore for all $\alpha\in\bbF_q$, we have $\eta(\alpha)=\alpha^{\nicefrac{(q-1)}{2}}$. The \emph{Legendre symbol}\index{Legendre symbol} $\Leg{\alpha}{q}$ for any $\alpha\in\bbF_q$ is defined by $\eta(\alpha)=\Leg{\alpha}{q}$
\end{note}

For any multiplicative character $\psi$ of $\bbF_q$, the value of $\psi(-1)$ is of special interest. We obviously have $\psi(-1)=\pm 1$. Let $\ord(\psi)=m$. Then $m\mid q-1$. Hence values of $\psi$ are $m^{th}$ roots of unity. Therefore $\psi(-1)$ can be $-1$ if $m$ is even. Let $g$ be a primitive element of $\bbF_q$ i.e. $\bbF_q^*$ is generated by $g$. Hence $\psi(g)=\zeta_m$ where $\zeta_m$ is the $m^{th}$ root of unity. If $m$ is even $$\psi(-1)=\psi\lt(g^{\nicefrac{(q-1)}2}\rt)=\zeta_m^{\nicefrac{(q-1)}2}$$Now $$\zeta_m^{\nicefrac{(q-1)}2}=-1 \iff \frac{q-1}2\equiv \frac{m}2\bmod m \iff \frac{q-1}m\equiv 1\bmod 2$$So $m$ should be even and $\frac{q-1}m$ should be odd. Therefore we have the following two observations.

\begin{observation}\label{obs:char-at--1}
  For any multiplicative character $\psi\in\Mq$, $\psi(-1)=-1$ if and only if $\frac{q-1}m$ is odd and $m$ is even.
\end{observation}
\begin{observation}
  If $\psi\in\Mq$ is the principal multiplicative character of $\bbF_q$ then $\psi(-1)=-1$. 
\end{observation}


Suppose $d\mid q-1$. Then for every multiplicative character $\psi$ of exponent $d$ and for every $\alpha\in \bbF_q^*$, we have $\psi(\alpha^d)=\psi^d(\alpha)=1$. Hence $\psi(\alpha)=1$ if $\alpha\in (\bbF_q^*)^{(d)}$. Conversely if $\psi(\alpha)=1$ for every $\alpha\in (\bbF_q^*)^{(d)}$ then $\psi^d=\psi_0$. Hence $\psi$ is a character of exponent $d$ then $\psi(\alpha)$ depends only on the cosets of $\quot{\bbF_q^*}{(\bbF_q^*)^{(d)}}$. Therefore there are precisely $d$ such characters of exponent $d$.

\begin{Theorem}{Character Sum over $d$th-Power Characters}{thm:char-power-sum-d}
  Suppose $d\mid q-1$. Suppose $\alpha\in\bbF_q^*$. Then $$\sum_{\psi\in\Mq^{(d)}}\chi(\alpha)=\begin{cases}
      d & \text{if }\alpha\in (\bbF_q^*)^{(d)} \\
      0 & \text{otherwise}
    \end{cases}$$
\end{Theorem}
\begin{proof}
  The characters of exponent $d$ are characters of the group $\quot{\bbF_q^*}{(\bbF_q^*)^{(d)}}$. Hence this follows from \obsref{obs:char-sum}.
\end{proof}

Again, the orthogonality relations in \obsref{obs:orthogonal-relation} when applied to multiplicative characters gives the following:
\begin{Theorem}{Orthogonality of Multiplicative Characters}{thm:mult-char-props}
  \begin{enumerate}[label=(\roman*)]
    \item For any multiplicative character $\psi$ of $\bbF_q$ $$\sum_{\alpha\in\bbF^*_q}\psi(\alpha)=\begin{cases}
              0   & \text{if $\psi\neq \psi_0$} \\
              q-1 & \text{if $\psi= \psi_0$}
            \end{cases}$$
    \item For any two multiplicative characters $\psi$ and $\lm$ of $\bbF_q$, $$\sum_{\alpha\in\bbF^*_q}\psi(\alpha)\cdot \ov{\lm(\alpha)}=\begin{cases}
              0   & \text{if $\psi\neq \lm$} \\
              q-1 & \text{if $\psi= \lm$}
            \end{cases}$$
    \item For any two elements $\alpha,\beta\in\bbF^*_q$, $$\sum_{\psi\in \Mq}\psi(\alpha)\cdot\ov{\psi(\beta)}=\begin{cases}
              0   & \text{if $\alpha\neq \beta$} \\
              q-1 & \text{if $\alpha=\beta$}
            \end{cases}$$
  \end{enumerate}
\end{Theorem}


Therefore the multiplicative characters form an orthonormal basis on all complex valued functions from $\bbF_q^*$ to $\bbC$. So we have the observation
\begin{observation}{obs:mult-char-fourier-basis}
  The set of multiplicative characters $\Mq$ of $\bbF_q$ forms an orthonormal basis for the vector space $W$ of all complex valued functions $\varphi\colon\bbF_q^*\to\bbC$.  
\end{observation}
Usually we extend the definition of multiplicative characters to the whole $\bbF_q$ so that we can define special character sums over the whole field. Let $\psi$ is a multiplicative character then $\psi(0)$ is defined to be $0$ if $\psi$ is nontrivial multiplicative character and $\psi(0)=1$ if $\psi$ is the trivial character.
\subsection{Lifting of Characters to Finite Extension}\label{sec:lifting-chars}%
\index{character!lifting}\index{lifting!of characters}%
Let $E$ be a finite extension field of $\bbF_q$. Let $\chi_1$ and $\mu_1$ be the canonical additive characters of $\bbF_q$ and $E$ respectively. Let $\tr_{E/\bbF_q}$ be the trace function from $E$ to $\bbF_q$.

Then $\chi_1$ and $\mu_1$ are related by the identity $$\chi_1(Tr_{E/\bbF_q}(\alpha))=\mu_1(\alpha) \quad \forall\ \alpha\in E$$Since the trace function satisfies transitivity property we have $\tr_E=\tr_{\bbF_q}\circ\tr_{E/\bbF_q}$. Hence for all $\alpha\in E$, $\mu_1(\alpha)=\chi_1(\tr_{E/\bbF_q}(\alpha))$.

Similarly if $\psi_1$ and $\lambda_1$ are principal multiplicative characters of $\bbF_q$ and $E$ respectively then, $$\lambda_1(\alpha)=\psi_1(\norm_{E/\bbF_q}(\alpha))\quad \forall\ \alpha\in E^*$$ where $\norm_{E/\bbF_q}$ is the \emph{norm function}\index{norm function}\index{norm!of a field extension} from $E$ to $\bbF_q$.

In the opposite way for any additive character $\chi$ of $\bbF_q$ or multiplicative character $\psi$ of $\bbF_q$ we can lift the characters to an additive and multiplicative character of $E$ respectively by taking $\chi\circ\tr_{E/\bbF_q}$ and $\psi\circ\norm_{E/\bbF_q}$ respectively. We denote these by $\chi^{(r)}\coloneqq \chi\circ\tr_{E/\bbF_q}$ and $\psi^{(r)}\coloneqq \psi\circ\norm_{E/\bbF_q}$ where $r=[E:\bbF_q]$.

In the following sections we will be working with hybrid character sums which are formed by taking product of multiple characters of both kind. Hence we will be using $\chi,\tau$ to denote additive characters and $\psi, \lambda$ to denote multiplicative characters.